%% ma: L10-B13-0089
%% lop: 10
%% chuong: 5
%% bai: 13
%% dang: 10.13.12
%% loai: TN4
%% muc: NB
%% dap_an: "A"
%% nguon: "chuyen-de-cac-so-dac-trung-cua-mau-so-lieu-khong-ghep-nhom-toan-10-chuong-trinh-moi.pdf (GV. Phan Nhật Linh) (P03), Bài 02 Dạng 2 Phần I Câu 15"
%% kien_thuc: "Lớp 10 Bài 13 (mốt)"
%% trang_thai: chinh-thuc
%% ghi_chu: "Tự duyệt theo yêu cầu giáo viên 22/10/2026."
\begin{ex}
Một cửa hàng giày nam thống kê cỡ giày bán được trong một tháng để biết nên nhập cỡ giày nào nhiều, kết quả như sau:
\begin{center}
\begin{tabular}{|c|c|c|c|c|c|c|c|}\hline
Cỡ giày & $37$ & $38$ & $39$ & $40$ & $41$ & $42$ & $43$\\\hline
Số lượng & $3$ & $5$ & $18$ & $21$ & $32$ & $28$ & $4$\\\hline
\end{tabular}
\end{center}
Căn cứ vào mẫu thống kê, cửa hàng nên nhập cỡ giày nào với số lượng nhiều nhất?
\choice{\True $41$}{$43$}{$38$}{$39$}
\loigiai{Mốt của mẫu số liệu là $41$ (tần số $32$ lớn nhất) nên nhập cỡ $41$ nhiều nhất. Chọn A.}
\end{ex}
